\chapter{Aplicación metodológica al caso de estudio}

Este capítulo aplica la metodología al proyecto de la Defensoría del Pueblo en San Cristóbal. Su función es mostrar cómo se construye la línea base, cómo se organizan los indicadores y cómo se comparan los escenarios de simulación. De esta forma, el capítulo funciona como el puente entre el método y la propuesta de optimización.

\section{Caracterización del edificio}

La caracterización describe las condiciones arquitectónicas, climáticas y constructivas del caso de estudio. Esta sección deberá incluir implantación, orientación, programa arquitectónico, sistema constructivo, materiales predominantes, configuración de envolvente, relación lleno-vacío, protecciones solares existentes y criterios de ventilación natural.

El objetivo de esta sección es reconocer las condiciones iniciales que inciden en el desempeño ambiental del proyecto. En el contexto de Galápagos, la lectura del edificio debe considerar no solo el consumo energético, sino también las restricciones de transporte, mantenimiento, disponibilidad material y fragilidad ambiental.

\section{Matriz de indicadores}

La matriz de indicadores organiza los criterios con los que se evaluará el edificio. A diferencia de una certificación completa, esta matriz funciona como una herramienta de comparación entre escenarios. Los indicadores se agrupan en cinco dimensiones: energía, confort térmico, iluminación natural, materiales y adaptación al contexto insular.

\begin{table}[H]
  \centering
  \caption{Dimensiones de evaluación del proyecto}
  \begin{tabular}{p{3.2cm}p{8.8cm}}
    \toprule
    Dimensión & Indicadores de análisis \\
    \midrule
    Energía & Demanda energética, cargas térmicas, consumo estimado y reducción porcentual frente a línea base. \\
    Confort térmico & Temperatura operativa, horas de confort, ventilación natural y riesgo de sobrecalentamiento. \\
    Iluminación natural & Aprovechamiento de luz natural, necesidad de iluminación artificial y control de deslumbramiento. \\
    Materiales & Huella ambiental, disponibilidad local o regional, durabilidad y comportamiento térmico. \\
    Contexto insular & Factibilidad logística, mantenimiento, impacto visual y pertinencia ecológica. \\
    \bottomrule
  \end{tabular}
\end{table}

\section{Construcción de la línea base}

La línea base corresponde al estado actual del proyecto antes de aplicar estrategias de optimización. Esta línea base debe documentar los parámetros usados en el modelo: geometría, orientación, materiales, valores térmicos, infiltración, ventilación, horarios de uso, cargas internas, sombreamiento y datos climáticos.

El resultado esperado de esta sección es una lectura clara de los puntos críticos del edificio. Entre ellos pueden aparecer ganancias solares excesivas, baja protección de la envolvente, ventilación insuficiente, alta demanda de enfriamiento o dependencia de sistemas activos.

\section{Escenarios de simulación}

Los escenarios de simulación se organizan para aislar el efecto de cada tipo de estrategia y luego evaluar su interacción. Esta estructura permite identificar si la mejora principal proviene del diseño pasivo, de la materialidad o de la combinación de ambos.

\begin{enumerate}
  \item Escenario base: reproduce las condiciones actuales del proyecto.
  \item Escenario pasivo: modifica variables de diseño pasivo, como control solar, ventilación natural, sombreamiento e inercia térmica.
  \item Escenario material: evalúa alternativas de materiales sostenibles y mejoras de envolvente.
  \item Escenario combinado: integra las mejores medidas pasivas y materiales para estimar el máximo potencial de optimización.
\end{enumerate}

\section{Análisis comparativo}

El análisis comparativo deberá presentar los resultados por indicador y por escenario. Para evitar una lectura dispersa, cada resultado deberá responder tres preguntas: qué cambia respecto a la línea base, cuánto cambia y qué implicación tiene para el proyecto.

\begin{table}[H]
  \centering
  \caption{Estructura de comparación entre escenarios}
  \begin{tabular}{p{3cm}p{2.5cm}p{2.5cm}p{2.5cm}p{2.5cm}}
    \toprule
    Indicador & Base & Pasivo & Material & Combinado \\
    \midrule
    Demanda energética & Completar & Completar & Completar & Completar \\
    Horas de confort & Completar & Completar & Completar & Completar \\
    Ganancias térmicas & Completar & Completar & Completar & Completar \\
    Desempeño lumínico & Completar & Completar & Completar & Completar \\
    Emisiones asociadas & Completar & Completar & Completar & Completar \\
    \bottomrule
  \end{tabular}
\end{table}

\section{Hallazgos para la optimización}

El cierre del capítulo sintetizará los hallazgos que deben pasar al capítulo siguiente. No se trata todavía de recomendar soluciones, sino de identificar evidencia: qué variable tuvo mayor incidencia, qué estrategia produjo mejor respuesta, qué medida resultó poco efectiva y qué restricciones condicionan la implementación.
